\Needspace{4\baselineskip}
\phantomsection
\label{app:rewrite-finite-existence}
\label{proof:rewrite-equilibrium-regimes}
\begin{proof}[Proof of Proposition~\ref{prop:rewrite-equilibrium-regimes}]

The proof has three tasks. First, Lemma~\ref{lem:rewrite-monopoly-choice}
gives a unique, smooth, globally optimal within-date service choice at given
states. I combine this choice with the other equilibrium equations to
identify candidate long-run values and verify their admissibility, including
positive consumption. Finding these values does not yet establish the
existence of a trajectory approaching them.

Second, I write the exact dynamic equations in coordinates with finite
candidate limits. The linearization identifies the stable directions; I
also check that they can accommodate every sufficiently nearby initial
state. Counting stable directions alone would not establish this property.
The stable-manifold and inverse-function theorems then select the initial
consumption and shadow price for those states. They establish convergent
solutions of the nonlinear equations for all future time, not merely
solutions of their linear approximation.

Third, I reconstruct the original variables and verify feasibility,
market clearing, and both transversality conditions. Household concavity
and Lemma~\ref{lem:rewrite-developer-verification} establish optimality.
For the developer, I verify the curvature conditions of
Lemma~\ref{lem:rewrite-hamiltonian-concavity} at every AI-efficiency level
reachable under alternative research policies from the initial efficiency.
The upper bound makes this possible for initial efficiency sufficiently
close to it. These final checks turn the constructed solutions of the
first-order conditions into equilibria. The result is local in initial
states, not in the time horizon.

\par\medskip\noindent\textit{Step 1. Feasible efficiency and candidate long-run values.}

Let \(D=\overline B-B\). The state equation with an upper bound gives
\[
 \frac{\dot D}{D}=-\frac{\chi}{\overline B}(BM)^\eta.
\]
Integration gives
\begin{equation}
 \overline B-B_t=(\overline B-B_0)
 \exp\left\{-\frac{\chi}{\overline B}
 \int_0^t(B_sM_s)^\eta\dd s\right\}>0.
 \label{eq:rewrite-cap-gap}
\end{equation}
Local integrability makes the exponent finite at every finite date, so
\(D_t>0\).

For the candidate limits, impose $B=\overline B$ only in the normalized
boundary equations. Actual trajectories will satisfy $B<\overline B$
at every finite date.


Write $\theta=(1-\alpha)/\alpha$. For $\sigma\ne1$, let
$\varphi=(\sigma-1)/\sigma$ and define
\[
 m(s)=1-\frac{1-s}{\sigma}-\alpha s,\qquad
 \mathcal H_\sigma(s)=
 \left[(1-\alpha)s\,m(s)
       \left(\frac{\omega_X}{s}\right)^{1/\varphi}\right]^\theta.
\]
Here $m(s)=1-e_X$ is the marginal-revenue factor, not research compute.
The static monopoly optimum has $m(s_X)>0$ by
Lemma~\ref{lem:rewrite-monopoly-choice}. The CES share, monopoly condition,
and production function, equations~\eqref{eq:rewrite-eq-sx},
\eqref{eq:rewrite-eq-monopoly}, and \eqref{eq:rewrite-eq-y}, imply
\begin{align}
 \frac{X}{AN}
 &=\left[\frac{\omega_Ls_X}{\omega_X(1-s_X)}\right]^{1/\varphi},
 &
 \frac{U}{Y}&=(1-\alpha)s_Xm(s_X), \label{eq:rewrite-frontier-static-shares}\\
 \frac{Y}{K}&=B^\theta\mathcal H_\sigma(s_X).
 \label{eq:rewrite-frontier-share-map}
\end{align}
The identities $Z=X(\omega_X/s_X)^{1/\varphi}$ and
$X=B(1-\alpha)s_Xm(s_X)Y$ give the output--capital ratio in
\eqref{eq:rewrite-frontier-share-map}.

If $0<\sigma<1$, positive marginal revenue requires
\[
 s>s_{\min}\equiv\frac{1-\sigma}{1-\alpha\sigma}.
\]
On $(s_{\min},1)$, both $m(s)$ and $s^{-1/(\sigma-1)}$ increase strictly.
Since
\[
 \mathcal H_\sigma(s)^{1/\theta}
 =(1-\alpha)\omega_X^{1/\varphi}
 m(s)s^{-1/(\sigma-1)},
\]
the share map increases from zero to
$\mathcal H_\sigma(1)=[(1-\alpha)^2\omega_X^{1/\varphi}]^\theta$.
Thus the equation
\begin{equation}
 \frac{\rho+\gamma+\delta}{\alpha}
 =\overline B^\theta\mathcal H_\sigma(\overline s_X)
 \label{eq:rewrite-terminal-share-root}
\end{equation}
has exactly one interior solution if and only if
$\overline B>\mathcal B(\sigma)$.

If $\sigma>1$, the share map decreases strictly on $(0,1)$.
To verify the sign, put $a=\sigma^{-1}-\alpha$, so $m(s)=\varphi+as$.
Its logarithmic derivative has the sign of $a(\sigma-2)s-\varphi$.
If $a(\sigma-2)\leq0$ this is negative. If $\sigma>2$ and $a>0$,
$a(\sigma-2)<(\sigma-2)/\sigma<\varphi$; if $1<\sigma<2$ and $a<0$,
$a(\sigma-2)<(1-\sigma^{-1})(2-\sigma)<\varphi$.
The map therefore decreases from infinity to $\mathcal H_\sigma(1)$.
Equation~\eqref{eq:rewrite-terminal-share-root} has exactly one interior
solution if and only if $\overline B<\mathcal B(\sigma)$.
Since $\mathcal R(B)=\alpha B^\theta\mathcal H_\sigma(1)-\delta$,
these calculations also derive
\eqref{eq:rewrite-ai-terminal-ratio}--\eqref{eq:rewrite-critical-frontier}.
The interior-share condition is therefore
$\mathcal R(\overline B)>\rho+\gamma$ for $\sigma<1$, but
$\mathcal R(\overline B)<\rho+\gamma$ for $\sigma>1$.

\Needspace{4\baselineskip}
For either interior-share case, let $x,y,u,k,c$ denote $X,Y,U,K,C$ divided
by $AN$, respectively. Construct the normalized terminal values sequentially:
\begin{align}
 \overline x&=
 \left[\frac{\omega_L\overline s_X}
 {\omega_X(1-\overline s_X)}\right]^{1/\varphi},
 &\overline u&=\frac{\overline x}{\overline B},\nonumber\\
 \overline y&=
 \frac{\overline u}{(1-\alpha)\overline s_Xm(\overline s_X)},
 &\overline k&=\frac{\alpha\overline y}{\rho+\gamma+\delta},\nonumber\\
 \overline c&=\overline y-\overline u
       -(\delta+n+\gamma)\overline k,
 &\overline r&=\rho+\gamma.
 \label{eq:rewrite-labor-terminal-values}
\end{align}
At $\sigma=1$, define $\beta=(1-\alpha)\omega_X$.
The exact static solution is
\[
 y=\left[(\beta^2B)^\beta k^\alpha\right]^{1/(1-\beta)},
 \qquad u=\beta^2y,\qquad x=Bu.
\]
The equation $y/k=(\rho+\gamma+\delta)/\alpha$ has exactly one positive
solution for every finite $\overline B>0$:
\begin{equation}
 \overline k=
 \left[(\beta^2\overline B)^\beta
 \left(\frac{\rho+\gamma+\delta}{\alpha}\right)^{\beta-1}
 \right]^{1/(1-\alpha-\beta)}.
 \label{eq:rewrite-unit-terminal-capital}
\end{equation}
Set
$\overline y=[(\rho+\gamma+\delta)/\alpha]\overline k$,
$\overline u=\beta^2\overline y$,
and $\overline x=\overline B\,\overline u$, with $\overline c,\overline r$
as in \eqref{eq:rewrite-labor-terminal-values}.

Consumption positivity does not require an additional restriction.
Positive marginal revenue implies $0<m(\overline s_X)<1$.
Since $0<\overline s_X<1$, in all three cases
\[
 \frac{\overline u}{\overline y}
 =(1-\alpha)\overline s_Xm(\overline s_X)<1-\alpha,\qquad
 \frac{(\delta+n+\gamma)\overline k}{\overline y}
 =\frac{\alpha(\delta+n+\gamma)}{\rho+\gamma+\delta}<\alpha.
\]
It follows that $\overline c/\overline y>0$.

\par\medskip\noindent\textit{Step 2. The three positive-labor-share cases.}

The cases in Proposition~\ref{prop:rewrite-capped-complements-existence},
\ref{prop:rewrite-capped-unit-existence}, and
\ref{prop:rewrite-capped-substitutes-labor-existence} share one
normalization. Divide capital, consumption, and the AI-efficiency shadow
price by effective labor, and multiply the remaining efficiency gap by
effective labor:
\begin{equation}
 k=\frac K{AN},\quad c=\frac C{AN},\quad
 d=(\overline B-B)AN,\quad \widetilde q=\frac q{AN},
 \quad \tau=(AN)^{-1}.
 \label{eq:rewrite-labor-proof-coordinates}
\end{equation}
Dividing by $AN$ removes effective-labor growth. Multiplying the shrinking
gap by $AN$ allows it to converge to a positive value:
$\dot d/d=n+\gamma-\dot B/(\overline B-B)$.
The variable $d$ measures the scaled gap, not its rate of closure.
The constants $\overline k,\overline d$ are the limiting values of $k,d$;
they are determined by the model, not additional parameters.
Also $\dot\tau=-(n+\gamma)\tau$, so $\tau\to0$.
These coordinates leave the model unchanged: $K=k/\tau$ and
$B=\overline B-d\tau$ at every finite date.

Given $k,B$, solve the unique static monopoly
problem for $y,x,u$, and set $r=\alpha y/k-\delta$.
Define $a=\dot B/(\overline B-B)$, the proportional rate of closure of the
AI-efficiency gap. The research condition gives
\[
 B=\overline B-d\tau,\qquad
 M=\left[\widetilde qd\,\chi\eta\frac{B^\eta}{\overline B}
       \right]^{1/(1-\eta)},\qquad
 a=\chi\frac{B^\eta M^\eta}{\overline B}.
\]
Substituting \eqref{eq:rewrite-eq-research} and
\eqref{eq:rewrite-equilibrium-dynamics} gives the exact dynamic equations
\begin{subequations}
\label{eq:rewrite-labor-proof-dynamics}
\begin{align}
 \frac{\dot k}{k}&=\frac{y-c-u-M\tau}{k}-\delta-n-\gamma,\\
 \frac{\dot c}{c}&=r-\rho-\gamma,\\
 \frac{\dot d}{d}&=n+\gamma-a,\\
 \frac{\dot{\widetilde q}}{\widetilde q}
 &=r-\frac{x}{\widetilde qB^2}
   -\eta\frac{ad\tau}{B}+a-n-\gamma,\\
 \dot\tau&=-(n+\gamma)\tau.
\end{align}
\end{subequations}
In particular, population growth cancels from the normalized consumption
equation; it is not subtracted a second time.

Complete the terminal construction by setting
\begin{equation}
 \overline M=
 \left[\frac{(n+\gamma)\overline B^{1-\eta}}{\chi}\right]^{1/\eta},
 \qquad
 \overline{\widetilde q}=\frac{\overline x}{\overline r\,\overline B^2},
 \qquad
 \overline d=\frac{\overline M}{\eta(n+\gamma)\overline{\widetilde q}}.
 \label{eq:rewrite-labor-terminal-research}
\end{equation}
All are positive. Direct substitution shows that
$(\overline k,\overline c,\overline d,\overline{\widetilde q},0)$
is a fixed point of \eqref{eq:rewrite-labor-proof-dynamics}.
This boundary point is a device for constructing paths with $\tau>0$ and
$B<\overline B$, not an admissible date-zero state itself.

I next verify saddle-path stability. Write $\mu=\eta/(1-\eta)$,
$r_k=\partial r/\partial\log k$, and
$a_k=\partial[(y-u)/k]/\partial\log k+\overline c/\overline k$,
evaluated at the fixed point. In
$(\log k,\log c,\log d,\log\widetilde q,\tau)$ coordinates the Jacobian has
two diagonal blocks: one for capital and consumption, and one for the scaled
efficiency gap and normalized shadow price,
\begin{equation}
 J_R=\begin{pmatrix}
 a_k&-\overline c/\overline k\\ r_k&0
 \end{pmatrix},
 \qquad
 J_B=\begin{pmatrix}
 -(n+\gamma)\mu&-(n+\gamma)\mu\\
 (n+\gamma)\mu&\overline r+(n+\gamma)\mu
 \end{pmatrix},
 \label{eq:rewrite-labor-proof-blocks}
\end{equation}
and the additional eigenvalue $-(n+\gamma)$. There can be forcing from
$\tau$ to both blocks and from $(k,c)$ to $(d,\widetilde q)$, but not in the
reverse directions at the fixed point.

To check $r_k<0$, let $s=s_X$, $u_\sigma=1/\sigma$, $e=e_X$, and let $D(s)$
be the positive curvature expression in
\eqref{eq:rewrite-monopoly-proof-curvature}. Implicit differentiation of the
monopoly condition gives
\[
 \frac{\partial\log X}{\partial\log K}
 =\frac{\alpha(1-e)}{D(s)},\qquad
 r_k=(r+\delta)(1-\alpha)
 \left[\frac{\alpha s(1-e)}{D(s)}-1\right].
\]
Moreover,
\[
 D(s)-\alpha s(1-e)
 =(1-s)\{u_\sigma(1-e)+
       (\alpha-u_\sigma)(1-u_\sigma)s\}>0.
\]
The sign is immediate when $u_\sigma\geq1$ or $u_\sigma\leq\alpha$.
For $\alpha<u_\sigma<1$, the braces are linear in $s$, with positive
endpoint values $u_\sigma(1-u_\sigma)$ and
$u_\sigma^2+\alpha(1-2u_\sigma)$.
The latter is positive if $u_\sigma\leq1/2$; otherwise $\alpha<u_\sigma$
implies it exceeds $u_\sigma(1-u_\sigma)>0$.
Thus $r_k<0$ in every interior-share case.

Both blocks have negative determinants:
\[
 \det J_R=(\overline c/\overline k)r_k<0,\qquad
 \det J_B=-(n+\gamma)\mu\overline r<0.
\]
There are therefore three stable and two unstable eigenvalues. The stable
direction in $J_R$ has a nonzero $k$ component, and the stable direction in
$J_B$ has a nonzero $d$ component. To verify that their number suffices to
select the jump variables, consider a vector in the stable subspace whose
state components $(k,d,\tau)$ are all zero. Its linearized $\tau$ path is
zero. The unforced capital--consumption block must then follow its stable direction;
its zero initial $k$ component forces its $c$ component to be zero as well.
With that forcing absent, the same argument in the research block forces
$\widetilde q=0$. The projection of the three-dimensional stable subspace
onto the three predetermined coordinates is therefore injective and hence
invertible. This argument also covers repeated stable eigenvalues.

Lemma~\ref{lem:rewrite-monopoly-choice} makes the static solution smooth.
The normalized vector field extends smoothly through $\tau=0$, since
$B=\overline B-d\tau$ remains positive nearby. The stable-manifold theorem
and the inverse-function theorem therefore give a local graph
$(c,\widetilde q)$ over $(k,d,\tau)$; see \citet{dyatlov2018}.
Choose a small neighborhood of the full fixed point and then a smaller
initial-state neighborhood $\mathcal U_L$ of
$(\overline k,\overline d,0)$. The size of this neighborhood depends on the parameters
and the case under consideration. In original initial stocks the restriction is
\begin{equation}
 (k_0,d_0,\tau_0)=
 \left(\frac{K_0}{A_0N_0},\,
       (\overline B-B_0)A_0N_0,\,
       \frac{1}{A_0N_0}\right)\in\mathcal U_L.
 \label{eq:rewrite-labor-local-initial}
\end{equation}
Closeness is required in all three coordinates; $B_0$ close to
$\overline B$ alone is not sufficient. The point with $\tau=0$
is a long-run limit, not a finite-date initial state.
Every physical initial state satisfying
\eqref{eq:rewrite-labor-local-initial} has a solution on this
graph that stays in the full neighborhood for all future time and converges
to the fixed point. It is unique among solutions with those properties.
Positivity of $\tau,d,c,\widetilde q,k$ and the original RSI technology
ensure $0<B_t<\overline B$ and $M_t>0$ at every finite date.
Reconstruction gives the original initial stocks, the equilibrium equations,
and the limiting rates in
\eqref{eq:rewrite-complements-terminal}--\eqref{eq:rewrite-substitutes-labor-terminal}.
Indeed, the normalized vector field tends to zero at its fixed point, so
the time derivatives of its smooth static ratios tend to zero as well.
This step, not convergence of ratios alone, establishes convergence of
the growth rates.
Optimality and transversality are verified below.

\par\medskip\noindent\textit{Step 3. The AI-dominated case.}

Let $\sigma>1$ and $\overline B>\mathcal B(\sigma)$.
Here $\varphi=(\sigma-1)/\sigma\in(0,1)$.
I seek an equilibrium in which capital grows faster than effective labor.
Along such a trajectory, $K/(AN)$ would diverge, so I normalize consumption
and the shadow value by capital and scale the remaining AI-efficiency gap
by capital:
\[
 h=(AN/K)^\varphi,\quad d=(\overline B-B)K,\quad
 c=C/K,\quad v=q/K,\quad \tau=(AN)^{-1}.
\]
Within this construction, $y,x,u$ denote $Y/K,X/K,U/K$, respectively.
The constant $\overline d$ now denotes the limit of $(\overline B-B)K$,
not of $(\overline B-B)AN$; no original state or control is redefined.
Then
\[
 B=\overline B-d\tau h^{1/\varphi},\quad
 M=\left[vd\,\chi\eta B^\eta/\overline B\right]^{1/(1-\eta)},\quad
 a=\chi B^\eta M^\eta/\overline B,\quad r=\alpha y-\delta.
\]
Capital's growth rate is $g_K=y-c-u-M\tau h^{1/\varphi}-\delta$.
The exact equations are
\begin{subequations}
\label{eq:rewrite-ai-proof-dynamics}
\begin{align}
 \dot h&=\varphi h(n+\gamma-g_K),\\
 \dot d/d&=g_K-a,\\
 \dot c/c&=n+r-\rho-g_K,\\
 \dot v/v&=r-x/(vB^2)-\eta ad\tau h^{1/\varphi}/B+a-g_K,\\
 \dot\tau&=-(n+\gamma)\tau.
\end{align}
\end{subequations}
The CES term divided by capital is
$[\omega_Lh+\omega_Xx^\varphi]^{1/\varphi}$.
At $h=0$ the monopoly equation has the strictly positive curvature
$D(1)=\alpha(1-\alpha)$. Its implicit static solution is therefore smooth
in $(h,B)$ near $(0,\overline B)$. For the remaining terms, replacing
$h^{1/\varphi}$ by $\max\{h,0\}^{1/\varphi}$ outside the physical domain
gives a $C^1$ extension because $1/\varphi>1$. The $C^1$ stable-manifold
theorem applies; see \citet[p.~7]{feldman2021stable}. The extension is only a
mathematical device and does not change any equation on $h\geq0$.

The limiting interest and aggregate growth rates are
\begin{equation}
 \overline r=\mathcal R(\overline B),\qquad
 \overline g=n+\overline r-\rho>n+\gamma.
 \label{eq:rewrite-substitutes-terminal-growth}
\end{equation}
At $h=\tau=0$, the terminal constants are
\begin{align}
 \overline y&=\frac{\mathcal R(\overline B)+\delta}{\alpha},&
 \overline u&=(1-\alpha)^2\overline y,&
 \overline x&=\overline B\,\overline u,\nonumber\\
 \overline r&=\mathcal R(\overline B),&
 \overline g&=n+\overline r-\rho,&
 \overline c&=\alpha(1-\alpha)\overline y+\rho-n>0,\nonumber\\
 \overline v&=\overline x/(\overline r\,\overline B^2),&
 \overline M&=
 \left[\overline g\,\overline B^{1-\eta}/\chi\right]^{1/\eta},&
 \overline d&=\overline M/(\eta\overline g\,\overline v).
 \label{eq:rewrite-ai-terminal-values}
\end{align}
The threshold implies $\overline g>n+\gamma>0$ and
$\overline r>\rho+\gamma>0$, so all constants are positive.
Substitution verifies the fixed point
$(0,\overline d,\overline c,\overline v,0)$.

The $h$ and $\tau$ eigenvalues are
$\varphi(n+\gamma-\overline g)<0$ and $-(n+\gamma)<0$.
The first equals $\varphi[\rho+\gamma-\mathcal R(\overline B)]$.
It becomes zero at $\mathcal R(\overline B)=\rho+\gamma$, which is why
this stable-manifold argument does not cover the threshold.
The consumption-ratio eigenvalue is $\overline c>0$.
The remaining block in $(\log d,\log v)$ is
\[
 \begin{pmatrix}
 -\overline g\mu&-\overline g\mu\\
 \overline g\mu&\overline r+\overline g\mu
 \end{pmatrix},
 \qquad \det=-\overline g\mu\overline r<0.
\]
The linearized system is triangular in the order
$(h,\tau),c,(d,v)$. A stable vector with zero state components $(h,d,\tau)$
has zero $h,\tau$ paths; its consumption component must be zero because the
unforced consumption equation is unstable. The research block then has zero
$d$ component and hence zero $v$ component, as in the previous construction.
Thus the three-dimensional stable subspace projects invertibly onto
$(h,d,\tau)$. The same graph argument gives a neighborhood
$\mathcal U_A$ of $(0,\overline d,0)$ such that the initial-state condition is
\begin{equation}
 \left[\left(\frac{A_0N_0}{K_0}\right)^\varphi,\;
       (\overline B-B_0)K_0,\; \frac{1}{A_0N_0}\right]\in\mathcal U_A.
 \label{eq:rewrite-ai-local-initial}
\end{equation}
It constructs a convergent trajectory for every admissible initial state
in this neighborhood, with a unique nearby convergent choice
of $C_0,q_0$ and positive quantities at every finite date.

The normalized vector field tends to zero, so the time derivatives of the
smooth positive ratios $Y/K$, $X/K$, and $C/K$ tend to zero.
Consequently $g_Y,g_X,g_C,g_K\to\overline g$ and $s_X\to1$.
Log differentiation of the CES share in
equation~\eqref{eq:rewrite-eq-sx} and the wage equation
\eqref{eq:rewrite-eq-w} gives the exact identity
\[
 g_w=g_Y-n+\frac{\sigma-1}{\sigma}s_X(n+\gamma-g_X).
\]
Taking limits yields
\[
 g_w\to\overline g-n+\frac{\sigma-1}{\sigma}(n+\gamma-\overline g)
 =\gamma+\frac{\overline r-\rho-\gamma}{\sigma}.
\]
This proves the growth and distribution limits in
Proposition~\ref{prop:rewrite-capped-substitutes-existence}.

\par\medskip\noindent\textit{Step 4. Optimality and transversality.}

In either construction let $\overline g$ denote the terminal aggregate
growth rate: $n+\gamma$ in the positive labor-share case and
$n+\overline r-\rho$ in the AI-dominated case. The constructions imply
\[
 a\to\overline g,\qquad
 M\to\left[\frac{\overline g\,\overline B^{1-\eta}}{\chi}\right]^{1/\eta}>0,
 \qquad g_q\to\overline g,\qquad g_B\to0,\qquad M/Y\to0.
\]
Both transversality expressions in Definition~\ref{def:rewrite-equilibrium}
have limiting logarithmic growth
\begin{equation}
 -\rho+n+g_K-g_C\to n-\rho<0,\qquad
 -r+g_q+g_B\to\overline g-\overline r=n-\rho<0.
 \label{eq:rewrite-finite-existence-tvcs}
\end{equation}
They therefore converge to zero, rather than merely remaining bounded.

Household utility is finite because $\log(C/N)$ grows at most linearly and
$\rho>n$, as imposed in the household problem. Log utility is concave, and the
household budget is linear at given prices and distributed profits. More explicitly,
let $\lambda_t=e^{-\rho t}N_t/C_t$, so $\dot\lambda_t=-r_t\lambda_t$ by
\eqref{eq:rewrite-euler}. Concavity of log utility and equal initial capital
give, for any admissible alternative $(\widetilde C,\widetilde K)$,
\[
 \int_0^T e^{-\rho t}N_t
 \log\frac{\widetilde C_t}{C_t}\,\dd t
 \leq\int_0^T\lambda_t(\widetilde C_t-C_t)\,\dd t
 =-\lambda_T(\widetilde K_T-K_T)\leq\lambda_TK_T\to0.
\]
Here $\widetilde K_T\geq0$ and the last limit is the household TVC.
Final-good optimality follows from concavity and constant returns.
The pointwise monopoly optimum is global by
Lemma~\ref{lem:rewrite-monopoly-choice}.

For the dynamic developer, feasibility and first-order conditions still
require a sufficiency argument. Let $S_t= A_tN_t$ in the positive-share
construction and $S_t=K_t$ in the AI-dominated construction, and let
$D_t=\exp(-\int_0^t r_s\,\dd s)$. Homogeneity of the static problem and
compactness of the normalized state neighborhood give a constant $C_\pi$
such that $0\leq\pi_t(b)\leq C_\pi S_t$ for every date and every
$b\in[B_0,\overline B]$. Since
$\mathrm d\log(D_tS_t)/\mathrm dt\to\overline g-\overline r=n-\rho<0$,
$\int_0^\infty D_tS_t\,\dd t<\infty$. Thus discounted operating profit
is uniformly bounded above across all alternative research policies.
The candidate's present value of net profit is finite, because its research compute is
bounded and $r_t\to\overline r>0$. An alternative with infinite discounted
research expenditure has payoff $-\infty$ and cannot improve on it.

It remains to verify the curvature conditions of
Lemma~\ref{lem:rewrite-hamiltonian-concavity},
\eqref{eq:rewrite-capped-concavity-gate} and
\eqref{eq:rewrite-depth-research-concavity} for every date and every alternative
AI-efficiency level in $[B_0,\overline B)$, not just on the selected path.
First fix a compact normalized-state neighborhood and an interval of
alternative AI-efficiency levels $[b_*,\overline B]$, with $0<b_*<\overline B$.
In the positive labor-share construction, $K/(AN)$ stays in a compact
neighborhood of $\overline k>0$. Smoothness of the static monopoly solution
bounds $\varepsilon_{B,t}(b)$ uniformly for these states and for AI-efficiency levels near
$\overline B$. In the AI-dominated construction, $h$ stays near zero, and
the same uniform bound follows from the smooth static limit
$\varepsilon_{B,t}(b)\to1/\alpha$. Thus in either case there is a finite constant
$E$ with $\varepsilon_{B,t}(b)-2\leq E$ throughout that state domain.
This bound is chosen before shrinking the initial-state neighborhood.
Because $\mu=\eta/(1-\eta)$ is finite for each $0<\eta<1$, choose
$b_{**}\in[b_*,\overline B)$ sufficiently close to $\overline B$ that
\[
 (1-b_{**}/\overline B)\max\{E,\mu-1,0\}<b_{**}/\overline B.
\]
Shrink the open neighborhood $\mathcal U_L$ or $\mathcal U_A$ so that
$B_0>b_{**}$.
This is possible because $B_0=\overline B-d_0\tau_0$ in the first
construction and $B_0=\overline B-d_0\tau_0h_0^{1/\varphi}$ in the second.
Every feasible research policy has $b\geq B_0$, so the two curvature
conditions hold throughout the entire alternative-state domain at every
date. Lemma~\ref{lem:rewrite-hamiltonian-concavity} therefore gives the
support inequality~\eqref{eq:rewrite-hamiltonian-support}.
Lemma~\ref{lem:rewrite-developer-verification} and the verified TVC
then establish global developer optimality without imposing
$\eta\leq1/2$. The curvature conditions follow from the primitives and the
chosen initial-state neighborhood.

All quantities remain finite at every finite date, and market clearing and
the original initial conditions hold by reconstruction. The paths are
therefore equilibria, not merely solutions of the
canonical equations. The result is local only in its set of initial
stocks. For $\sigma>1$, at $\overline B=\mathcal B(\sigma)$ the interior-share terminal
point reaches the boundary and a stable eigenvalue in the boundary
normalization becomes zero; this proof does not apply. For
$\sigma<1$ and $\overline B\leq\mathcal B(\sigma)$, it does not rule out
equilibria with different long-run ratios. Nor does it imply existence at
arbitrary initial stocks or in the uncapped limit.

\par\medskip\noindent\textit{Step 5. Initial-state sets and the four conclusions.}

The preimage of each open neighborhood under its continuous
initial-state map, restricted to $A_0,N_0,K_0>0$ and
$0<B_0<\overline B$, is open in the original initial stocks.
It is nonempty. In the positive-labor-share cases, choose
$k_0,d_0$ close to $\overline k,\overline d$ and $\tau_0>0$
sufficiently small. Then $A_0N_0=1/\tau_0$,
$K_0=k_0/\tau_0$, and $B_0=\overline B-d_0\tau_0$ are admissible.
In the AI-dominated case, choose $d_0$ near $\overline d$ and
$h_0,\tau_0>0$ sufficiently small. The inverse relations are
$A_0N_0=1/\tau_0$, $K_0=1/(\tau_0h_0^{1/\varphi})$, and
$B_0=\overline B-d_0\tau_0h_0^{1/\varphi}$.
Any positive factorization of $A_0N_0$ supplies $A_0,N_0$.
These choices also satisfy the strict efficiency restriction in Step 4.

For each initial state in the resulting set, the local graph selects
$C_0,q_0>0$. These jump values are unique among the trajectories
that remain near the indicated normalized limit and converge to it.
The argument does not assert global uniqueness.

For $\sigma<1$ with $\overline B>\mathcal B(\sigma)$, and for
$\sigma>1$ with $\overline B<\mathcal B(\sigma)$, Step 1 gives a unique
interior $\overline s_X$. Step 2 then yields
$g_Y-n\to\gamma$, $g_w\to\gamma$, and $r\to\rho+\gamma$.
At $\sigma=1$, the exact CES limit gives $s_X=\omega_X$ at every date,
and the same construction works for every finite positive upper bound.
These establish the first three cases of the proposition.
Step 3 gives $s_X\to1$, $r\to\mathcal R(\overline B)$,
$g_Y-n\to\mathcal R(\overline B)-\rho$, and the wage-growth limit in
\eqref{eq:rewrite-substitutes-wage-growth}, establishing the fourth case.
Step 4 verifies the agent-optimality and transversality requirements of
Definition~\ref{def:rewrite-equilibrium} for all these trajectories.
\end{proof}

%\Needspace{20\baselineskip}
%\begin{proposition}[Wage growth and labor's income share]
%\label{prop:rewrite-output-wage-wedge}
%Consider the long-run equilibria constructed in Proposition~\ref{prop:rewrite-equilibrium-regimes}. Then
%\[
% \lim_{t\to\infty}g_w>\gamma
% \quad\Longleftrightarrow\quad
% \lim_{t\to\infty}\frac{wL}{Y}=0,
%\]
%and either condition holds precisely when the equilibrium is AI-dominated. Along such an equilibrium, output per person grows asymptotically faster than the real wage:
%\begin{equation}
% \lim_{t\to\infty}\bigl[(g_Y-n)-g_w\bigr]
% =\left(1-\frac{1}{\sigma}\right)
%  (\overline r-\rho-\gamma)>0.
% \label{eq:rewrite-output-wage-wedge}
%\end{equation}
%Because $L=N$, the same wedge is the asymptotic rate at which labor's income share falls:
%\begin{equation}
% \lim_{t\to\infty}\frac{\mathrm d}{\mathrm dt}
% \log\left(\frac{wL}{Y}\right)
% =-\left(1-\frac{1}{\sigma}\right)
%  (\overline r-\rho-\gamma)<0.
% \label{eq:rewrite-labor-share-decline-rate}
%\end{equation}
%Equivalently,
%\begin{equation}
% -\lim_{t\to\infty}\frac{\mathrm d}{\mathrm dt}
%  \log\left(\frac{wL}{Y}\right)
% =(\sigma-1)\lim_{t\to\infty}(g_w-\gamma)>0.
% \label{eq:rewrite-wage-premium-share-decline}
%\end{equation}
%\end{proposition}

%\Needspace{4\baselineskip}
%\phantomsection
%\label{proof:rewrite-output-wage-wedge}
%\begin{proof}[Proof of Proposition~\ref{prop:rewrite-output-wage-wedge}]
%The first three cases of Proposition~\ref{prop:rewrite-equilibrium-regimes} give $g_w\to\gamma$ and a positive limiting labor share, using $wL/Y=(1-\alpha)(1-s_X)$. Proposition~\ref{prop:rewrite-capped-substitutes-existence} implies both $g_w\to\gamma+(\overline r-\rho-\gamma)/\sigma>\gamma$ and $wL/Y\to0$. This proves the equivalence. Subtract the wage-growth limit in \eqref{eq:rewrite-substitutes-wage-growth} from the per-person output-growth limit in \eqref{eq:rewrite-substitutes-ai-limits}. Moreover, $\mathrm d\log(wL/Y)/\mathrm dt=g_w+n-g_Y$ because $L=N$. Finally, combine \eqref{eq:rewrite-labor-share-decline-rate} with \eqref{eq:rewrite-substitutes-wage-growth}.
%\end{proof}
